\section*{Sets and Indices}

No nontrivial sets or indices are used; this is a single-period continuous production planning model with aggregate process-capacity constraints.

\section*{Parameters}

\[
t^{A}_{1}=2
\]
time required in process 1 per unit of product A (code: \texttt{a\_p1}).

\[
t^{A}_{2}=3
\]
time required in process 2 per unit of product A (code: \texttt{a\_p2}).

\[
t^{B}_{1}=3
\]
time required in process 1 per unit of product B (code: \texttt{b\_p1}).

\[
t^{B}_{2}=4
\]
time required in process 2 per unit of product B (code: \texttt{b\_p2}).

\[
T_{1}=16
\]
total available time in process 1 (code: \texttt{T1}).

\[
T_{2}=24
\]
total available time in process 2 (code: \texttt{T2}).

\[
\gamma = 2
\]
units of by-product C generated per unit of product B (code: \texttt{c\_per\_b}).

\[
\overline{C}^{\mathrm{sell}}=5
\]
maximum amount of by-product C that can be sold (code: \texttt{max\_c\_sales}).

\[
\pi_A = 4
\]
profit per unit of product A (code: \texttt{profit\_a}).

\[
\pi_B = 10
\]
profit per unit of product B (code: \texttt{profit\_b}).

\[
\pi_C = 3
\]
profit per unit of by-product C sold (code: \texttt{profit\_c}).

\[
d_C = 2
\]
disposal cost per unit of by-product C disposed (code: \texttt{disposal\_cost}).

\section*{Decision Variables}

\[
x_A \ge 0
\]
units of product A produced (code: \texttt{xA}).

\[
x_B \ge 0
\]
units of product B produced (code: \texttt{xB}).

\[
c^{\mathrm{sell}} \ge 0
\]
units of by-product C sold (code: \texttt{cSold}).

\[
c^{\mathrm{disp}} \ge 0
\]
units of by-product C disposed (code: \texttt{cDisposed}).

\section*{Objective}

Maximize total profit:
\[
\max \;
\pi_A x_A + \pi_B x_B + \pi_C c^{\mathrm{sell}} - d_C c^{\mathrm{disp}}.
\]

With the instance data, this is
\[
\max \; 4x_A + 10x_B + 3c^{\mathrm{sell}} - 2c^{\mathrm{disp}}.
\]

\section*{Constraints}

Process 1 capacity:
\[
t^{A}_{1}x_A + t^{B}_{1}x_B \le T_1.
\]

Process 2 capacity:
\[
t^{A}_{2}x_A + t^{B}_{2}x_B \le T_2.
\]

By-product C balance:
\[
c^{\mathrm{sell}} + c^{\mathrm{disp}} = \gamma x_B.
\]

Upper bound on sales of by-product C:
\[
c^{\mathrm{sell}} \le \overline{C}^{\mathrm{sell}}.
\]

Nonnegativity:
\[
x_A,\;x_B,\;c^{\mathrm{sell}},\;c^{\mathrm{disp}} \ge 0.
\]

Substituting the numerical values:
\[
2x_A + 3x_B \le 16,
\]
\[
3x_A + 4x_B \le 24,
\]
\[
c^{\mathrm{sell}} + c^{\mathrm{disp}} = 2x_B,
\]
\[
c^{\mathrm{sell}} \le 5,
\]
\[
x_A,\;x_B,\;c^{\mathrm{sell}},\;c^{\mathrm{disp}} \ge 0.
\]

\section*{Notes on Implementation}

The implementation in \texttt{current\_heuristic.py} builds exactly the above continuous linear program using PuLP-compatible syntax and solves it with \texttt{GUROBI\_CMD}. No integrality restrictions are imposed, so all variables are continuous.

The model does not introduce an explicit production variable for by-product C; instead, C generation is enforced solely through the balance equation
\[
c^{\mathrm{sell}} + c^{\mathrm{disp}} = \gamma x_B.
\]
Thus every unit of C generated from product B must be either sold or disposed.

The formulation should not be simplified further when documenting the implemented model; in particular, \texttt{cSold} and \texttt{cDisposed} remain explicit decision variables rather than being eliminated algebraically.
